% logo_utn.tex — Símbolo de la UTN («la Araña», Ricardo Blanco) en TikZ, vectorial.
%
% Geometría tomada del vector oficial del manual de identidad de la UTN
% (identidad.frc.utn.edu.ar, «Manual de uso UTN FRC», pág. 6, «Retícula
% constructiva»). Unidad = ancho del trazo:
%   - barra vertical de 1 de ancho, alto total 5,8963;
%   - barra horizontal de 5 × 1, centrada;
%   - arriba, media corona circular (radios 2,5 y 1,5) con centro en el borde
%     superior, abierta hacia arriba; abajo, la misma reflejada.
%
% Uso: \logoutn{<alto>}   (ej. \logoutn{1.9cm}); color = color del texto actual.
% Lo usan la carátula del PDF de lab_fisica (lab-fisica.lua) y logo_utn_png.tex
% (PNG para Word, generado una sola vez por comun.mk en logos/build/).

\providecommand{\logoutn}[1]{%
  \begingroup
  \pgfmathsetmacro{\logoescala}{#1/5.8963}%
  \begin{tikzpicture}[x=\logoescala pt, y=\logoescala pt, baseline=(current bounding box.south)]
    \def\h{2.94815}
    \fill (-0.5,-\h) rectangle (0.5,\h);
    \fill (-2.5,-0.5) rectangle (2.5,0.5);
    \fill (-2.5,\h) arc[start angle=180, end angle=360, radius=2.5]
          -- (1.5,\h) arc[start angle=0, end angle=-180, radius=1.5] -- cycle;
    \fill (-2.5,-\h) arc[start angle=180, end angle=0, radius=2.5]
          -- (1.5,-\h) arc[start angle=0, end angle=180, radius=1.5] -- cycle;
  \end{tikzpicture}%
  \endgroup
}
